% BHU PYQ -- Transcribed & Corrected
% Paper: CHEMJ31: Inorganic Chemistry-I
% Exam: B.Sc. (Hons) Semester III Examination, 2025-26 (NEP)
% Ref Code: CECONF/N/2025-26/056 (A)
% Category: Major
% Department: Chemistry

\documentclass[11pt,a4paper]{article}
\usepackage[a4paper,top=2.5cm,bottom=2.5cm,left=2.8cm,right=2.8cm,headheight=14pt]{geometry}
\usepackage{amsmath,amssymb}
\usepackage[T1]{fontenc}
\usepackage[utf8]{inputenc}
\usepackage{lmodern,microtype,fancyhdr,enumitem,hyperref,lastpage,parskip}

\hypersetup{
    colorlinks=true,
    urlcolor=black,
    linkcolor=black,
    pdftitle={CHEMJ31: Inorganic Chemistry-I -- BHU Sem III 2025-26 (NEP)}
}

\pagestyle{fancy}
\fancyhf{}
\renewcommand{\headrulewidth}{0.4pt}
\renewcommand{\footrulewidth}{0.4pt}
\fancyhead[L]{\small \textit{Banaras Hindu University}}
\fancyhead[R]{\small \textit{CHEMJ31: Inorganic Chemistry-I (2025-26)}}
\fancyfoot[C]{\small Page \thepage\ of \pageref{LastPage}}

\newcommand{\pts}[1]{\hfill \textit{[#1]}}

\begin{document}

\begin{center}
  {\large\bfseries BANARAS HINDU UNIVERSITY}\\[2pt]
  {\normalsize B.Sc. (Hons) Semester III Examination, 2025-26 (NEP)}\\[6pt]
  \rule{0.8\linewidth}{0.5pt}\\[6pt]
  {\large\bfseries CHEMISTRY}\\[4pt]
  {\large\bfseries Paper Code: CHEMJ31 : Inorganic Chemistry-I}\\[2pt]
  {\small Ref. No.: CECONF/N/2025-26/056 (A)}\\[4pt]
  {\normalsize Time: 2$\frac{1}{2}$ Hours\hfill Maximum Marks: 70}
\end{center}

\rule{\linewidth}{0.4pt}

\smallskip

\noindent \textit{Instructions: Answer five questions including Question 1 which is compulsory. The figures in the right-hand margin indicate marks.}

\bigskip

\noindent \textbf{Question 1.} Answer all parts:
\begin{enumerate}[label=(\alph*)]
  \item Compare the products of reaction when normal-, per- and super oxides of sodium are reacted separately with water. \pts{1$\frac{1}{2}$}
  \item Why do beryllium salts rarely have more than four molecules of water of crystallization associated with the metal ion? \pts{1}
  \item Write the correct formulation of $\text{GaCl}_2$ indicating the oxidation state of Ga. \pts{1$\frac{1}{2}$}
  \item Write the name and formula of the most widely used nitrogenous fertilizer along with one chemical reaction for its synthesis. \pts{2}
  \item ``The melting and boiling points for $\text{NH}_3$ are out of line as compared to those of its heavier congeners.'' Explain, why. \pts{2}
  \item Write the chemical composition for Portland cement. Why is slow setting of the same desired? How is the same achieved? \pts{2}
  \item Comment upon the stability of singlet and triplet $\text{O}_2$ along with mentioning of their energy values in kJ/mole. What happens when the singlet $\text{O}_2$ is allowed to react with 1,3-buta-diene? \pts{2}
  \item Write the names of $\text{OF}_2$ and $\text{Cl}_2\text{O}$ and draw their structures mentioning the hybridization state of the central atom in each case. Why is the Cl-O-Cl bond angle greater than that of F-O-F? \pts{2}
\end{enumerate}

\bigskip

\noindent \textbf{Question 2.}
\begin{enumerate}[label=(\alph*)]
  \item What do you mean by boranes? Illustrate 2e-2c and 2e-3c kinds of bonds in one lower and one higher borane. Write general molecular formulae for closo-, nido-, arachno- and hypho-boranes. \pts{4}
  \item Why are Group 1 elements (i) univalent (ii) largely ionic (iii) poor complexing agents? \pts{2}
  \item Compare multiple bonding in $\text{C}_2\text{H}_4$ and $\text{P}_4\text{O}_{10}$ giving suitable sketches. \pts{2}
  \item Why is $\text{CO}_2$ a gas while $\text{SiO}_2$ a solid at room temperature? \pts{2}
  \item Describe the structure of $\text{SO}_3$ in its solid state. \pts{2}
  \item Write a balanced chemical reaction showing preparation of $\text{H}_2\text{O}_2$. Describe peculiar point of its structure also. \pts{2}
\end{enumerate}

\bigskip

\noindent \textbf{Question 3.}
\begin{enumerate}[label=(\alph*)]
  \item Comment upon the acidity of $\text{H}_3\text{BO}_3$ in aqueous medium. How can its acidity be modulated by (i) taking its higher concentration (ii) adding a compound having cis-diol? \pts{6}
  \item Draw the structures of $\text{C}_3\text{O}_2$ and $\text{C}_{12}\text{O}_9$. Give their preparations also. \pts{4}
  \item ``The dipole moment of $\text{NH}_3$ is much greater than that of $\text{NF}_3$.'' Explain, why. \pts{2}
  \item How is an electride type complex prepared from Cryptand-222. Comment in very brief upon its magnetic behavior also. \pts{2}
\end{enumerate}

\bigskip

\noindent \textbf{Question 4.}
\begin{enumerate}[label=(\alph*)]
  \item Give a detailed account of the classification of silicates along with neat sketches of each. \pts{6}
  \item Discuss in brief the following properties of hydrides of Group 15: (i) Ease of preparation (ii) Stability (iii) Reducing power (iv) Ability to act as electron pair donors. \pts{6}
  \item ``$\text{BH}_3$ undergoes dimerization while $\text{BF}_3$ does not.'' Explain, why. \pts{2}
\end{enumerate}

\bigskip

\noindent \textbf{Question 5.}
\begin{enumerate}[label=(\alph*)]
  \item Compare the structure of diamond and graphite. Comment upon their relative thermodynamic stability and conducting behavior also. Describe the preparation of fullerenes from graphite in very brief. \pts{6}
  \item Discuss the structure of various types of oxides of Nitrogen. Draw the structures of $\text{HNO}_2$ and $\text{HNO}_3$ and mention names along with formulae for their anhydrides. Give a brief account of preparation of $\text{HNO}_3$ through Ostwald's process. \pts{6}
  \item Write the chemical composition of high alumina cement. Discuss its superiority over Portland cement in very brief. \pts{2}
\end{enumerate}

\bigskip

\noindent \textbf{Question 6.}
\begin{enumerate}[label=(\alph*)]
  \item What do you mean by Zeolite? Write its general formula. How the same can act as molecular sieve? Support your answer giving a suitable example. \pts{4}
  \item Define glass and discuss its composition. Give a brief account of different kinds of additions for having glass of desired property. \pts{4}
  \item Describe the structures for various kinds of oxoacids of sulfur along with mentioning of oxidation state of sulfur atoms in each case. \pts{6}
\end{enumerate}

\bigskip

\noindent \textbf{Question 7.}
\begin{enumerate}[label=(\alph*)]
  \item Define fertilizer. What are their different types? \pts{4}
  \item ``HF is the weakest acid among all the hydra acids of halogen family in their aqueous state.'' Explain, why. \pts{2}
  \item Describe one method by which hydrogen peroxide is prepared. Give the structure of $\text{H}_2\text{O}_2$ in the gas phase and mention its peculiar feature. \pts{3}
  \item Give the names of four different types of oxoacids of chlorine along with their molecular and structural formulae. Comment upon their relative strength giving reasons. \pts{4}
  \item Why is $\text{CCl}_4$ unaffected by water while $\text{SiCl}_4$ is rapidly hydrolysed? \pts{1}
\end{enumerate}

\bigskip
\begin{center}
  \textbf{***}
\end{center}

\end{document}
